(1) 传播矩阵 $P(d)=\begin{pmatrix}1&d\\0&1\end{pmatrix}$，薄透镜矩阵 $L(f)=\begin{pmatrix}1&0\\-1/f&1\end{pmatrix}$。总矩阵 $M=L(f_2)P(d)L(f_1)$，相乘得 $A=1-d/f_1$、$B=d$、$C=-(1/f_1+1/f_2-d/(f_1f_2))$、$D=1-d/f_2$，并满足 $AD-BC=1$。
(2) 一般近轴系统有效焦距满足 $C=-1/f_{\rm eff}$，所以 $1/f_{\rm eff}=1/f_1+1/f_2-d/(f_1f_2)$；等价为 $f_{\rm eff}=f_1f_2/(f_1+f_2-d)$，前提为分母不为零。不能直接将两片分离透镜的焦度相加，漏掉间隔修正。
(3) 平行入射时起始角 $\theta_{\rm in}=0$，第二透镜后光线为 $y_{
m out}=Ay_{
m in}$、$\theta_{
m out}=Cy_{
m in}$。再自由传播 $s$ 到轴上，要求 $A+sC=0$，故后焦距 $s=-A/C=f_{\rm eff}(1-d/f_1)=f_2(f_1-d)/(f_1+f_2-d)$，可为负值表示虚焦点。
(4) 无焦条件是 $C=0$，即 $d=f_1+f_2$；此时平行束出射仍平行，系统有效焦距形式上无穷大。若 $d\to0$，恢复贴合薄透镜焦度相加 $1/f_{\rm eff}=1/f_1+1/f_2$；若 $f_2\to\infty$，第二片趋于无光焦度平板，$f_{\rm eff}\to f_1$，但后焦点相对于第二片位置变为 $f_1-d$。
